
\documentclass{article}
\usepackage{colortbl}
\usepackage{makecell}
\usepackage{multirow}
\usepackage{supertabular}

\begin{document}

\newcounter{utterance}

\centering \large Interaction Transcript for game `dond', experiment `coop\_en', episode 8 with refineonly{-}Qwen3.5{-}27B.
\vspace{24pt}

{ \footnotesize  \setcounter{utterance}{1}
\setlength{\tabcolsep}{0pt}
\begin{supertabular}{c@{$\;$}|p{.15\linewidth}@{}p{.15\linewidth}p{.15\linewidth}p{.15\linewidth}p{.15\linewidth}p{.15\linewidth}}
    \# & $\;$A & \multicolumn{4}{c}{Game Master} & $\;\:$B\\
    \hline

    \theutterance \stepcounter{utterance}  
    & & \multicolumn{4}{p{0.6\linewidth}}{
        \cellcolor[rgb]{0.9,0.9,0.9}{
            \makecell[{{p{\linewidth}}}]{
                \texttt{\tiny{[P1$\langle$GM]}}
                \texttt{You are playing a cooperative negotiation game in which you have to agree on how to divide a set of items among you and another player.} \\
\\ 
\texttt{Rules:} \\
\texttt{(a) You and the other player are given a set of items. Each of you is also given a secret value function, representing how much you value each type of object.} \\
\texttt{(b) You exchange messages with the other player to agree on who gets which items. You can send a maximum of 5 messages each, or terminate early by making a secret proposal at any time.} \\
\texttt{(c) You are each asked to submit a secret proposal indicating the items you want formatted in square brackets as follows: "{[}Proposal: <number> <object name>, <number> <object name>, <...>{]}"} \\
\texttt{(d) If your proposals are complementary, i.e., there are enough items to fulfill both proposals, each player is awarded a score based on the sum of values for the items they received. Otherwise, both of you get zero points.} \\
\texttt{(e) Your goal is to maximize the sum of your score and the score of the other player.} \\
\\ 
\texttt{Let us start.} \\
\\ 
\texttt{The set of available items is:} \\
\\ 
\texttt{2 potatoes, 1 glue, 1 berry, 3 carrots} \\
\\ 
\texttt{Your secret value function is:} \\
\\ 
\texttt{potato: 1, glue: 2, berry: 3, carrot: 1} \\
\\ 
\texttt{IMPORTANT: Your messages, unless it is the secret proposal, are directly transmitted to the other player, so do not include any response to the rules or text announcing your message. To make a secret proposal, use the indicated format. Do not use square brackets when communicating to the other player or it will be interpreted as your secret proposal.} \\
            }
        }
    }
    & \\ \\

    \theutterance \stepcounter{utterance}  
    & \multicolumn{4}{p{0.6\linewidth}}{
        \cellcolor[rgb]{0.95,0.95,0.95}{
            \makecell[{{p{\linewidth}}}]{
                \texttt{\tiny{[P2$\rightarrow$P1]}}
                \texttt{{[}DRAFT{]}} \\
\texttt{Hi! Let's figure out the best way to split these items to maximize our combined score. I have my own values, but I'm sure you have yours too.} \\
\\ 
\texttt{Here is what I have:} \\
\texttt{{-} 2 potatoes} \\
\texttt{{-} 1 glue} \\
\texttt{{-} 1 berry} \\
\texttt{{-} 3 carrots} \\
\\ 
\texttt{My values are: potato (1), glue (2), berry (3), carrot (1).} \\
\\ 
\texttt{Could you share your value function for these items? That way, we can see who values what the most and make sure the items go to the person who gets the most points from them.} \\
            }
        }
    }
    & & \\ \\

    \theutterance \stepcounter{utterance}  
    & \multicolumn{4}{p{0.6\linewidth}}{
        \cellcolor[rgb]{0.9,0.9,0.9}{
            \makecell[{{p{\linewidth}}}]{
                \texttt{\tiny{[P1$\rangle$GM]}}
                \texttt{Hi! Let's figure out the best way to split these items to maximize our combined score. I have my own values, but I'm sure you have yours too.} \\
\\ 
\texttt{Here is what I have:} \\
\texttt{{-} 2 potatoes} \\
\texttt{{-} 1 glue} \\
\texttt{{-} 1 berry} \\
\texttt{{-} 3 carrots} \\
\\ 
\texttt{My values are: potato (1), glue (2), berry (3), carrot (1).} \\
\\ 
\texttt{Could you share your value function for these items? That way, we can see who values what the most and make sure the items go to the person who gets the most points from them.} \\
            }
        }
    }
    & & \\ \\

    \theutterance \stepcounter{utterance}  
    & & & \multicolumn{2}{p{0.3\linewidth}}{
        \cellcolor[rgb]{0.9,0.9,0.9}{
            \makecell[{{p{\linewidth}}}]{
                \texttt{\tiny{[GM$|$GM]}}
                \texttt{continue} \\
            }
        }
    }
    & & \\ \\

    \theutterance \stepcounter{utterance}  
    & & & \multicolumn{2}{p{0.3\linewidth}}{
        \cellcolor[rgb]{0.9,0.9,0.9}{
            \makecell[{{p{\linewidth}}}]{
                \texttt{\tiny{[GM$|$GM]}}
                \texttt{Hi! Let's figure out the best way to split these items to maximize our combined score. I have my own values, but I'm sure you have yours too.} \\
\\ 
\texttt{Here is what I have:} \\
\texttt{{-} 2 potatoes} \\
\texttt{{-} 1 glue} \\
\texttt{{-} 1 berry} \\
\texttt{{-} 3 carrots} \\
\\ 
\texttt{My values are: potato (1), glue (2), berry (3), carrot (1).} \\
\\ 
\texttt{Could you share your value function for these items? That way, we can see who values what the most and make sure the items go to the person who gets the most points from them.} \\
            }
        }
    }
    & & \\ \\

    \theutterance \stepcounter{utterance}  
    & & \multicolumn{4}{p{0.6\linewidth}}{
        \cellcolor[rgb]{0.9,0.9,0.9}{
            \makecell[{{p{\linewidth}}}]{
                \texttt{\tiny{[GM$\rangle$P2]}}
                \texttt{You are playing a cooperative negotiation game in which you have to agree on how to divide a set of items among you and another player.} \\
\\ 
\texttt{Rules:} \\
\texttt{(a) You and the other player are given a set of items. Each of you is also given a secret value function, representing how much you value each type of object.} \\
\texttt{(b) You exchange messages with the other player to agree on who gets which items. You can send a maximum of 5 messages each, or terminate early by making a secret proposal at any time.} \\
\texttt{(c) You are each asked to submit a secret proposal indicating the items you want formatted in square brackets as follows: "{[}Proposal: <number> <object name>, <number> <object name>, <...>{]}"} \\
\texttt{(d) If your proposals are complementary, i.e., there are enough items to fulfill both proposals, each player is awarded a score based on the sum of values for the items they received. Otherwise, both of you get zero points.} \\
\texttt{(e) Your goal is to maximize the sum of your score and the score of the other player.} \\
\\ 
\texttt{Let us start.} \\
\\ 
\texttt{The set of available items is:} \\
\\ 
\texttt{2 potatoes, 1 glue, 1 berry, 3 carrots} \\
\\ 
\texttt{Your secret value function is:} \\
\\ 
\texttt{potato: 1, glue: 0, berry: 2, carrot: 2} \\
\\ 
\texttt{IMPORTANT: Your messages, unless it is the secret proposal, are directly transmitted to the other player, so do not include any response to the rules or text announcing your message. To make a secret proposal, use the indicated format. Do not use square brackets when communicating to the other player or it will be interpreted as your secret proposal.} \\
\\ 
\\ 
\texttt{Hi! Let's figure out the best way to split these items to maximize our combined score. I have my own values, but I'm sure you have yours too.} \\
\\ 
\texttt{Here is what I have:} \\
\texttt{{-} 2 potatoes} \\
\texttt{{-} 1 glue} \\
\texttt{{-} 1 berry} \\
\texttt{{-} 3 carrots} \\
\\ 
\texttt{My values are: potato (1), glue (2), berry (3), carrot (1).} \\
\\ 
\texttt{Could you share your value function for these items? That way, we can see who values what the most and make sure the items go to the person who gets the most points from them.} \\
            }
        }
    }
    & \\ \\

    \theutterance \stepcounter{utterance}  
    & & & \multicolumn{4}{p{0.6\linewidth}}{
        \cellcolor[rgb]{0.95,0.95,0.95}{
            \makecell[{{p{\linewidth}}}]{
                \texttt{\tiny{[P1$\rightarrow$P2]}}
                \texttt{{[}DRAFT{]}} \\
\texttt{Hi! Thanks for sharing your values. Here is my value function:} \\
\\ 
\texttt{{-} potato: 1} \\
\texttt{{-} glue: 0} \\
\texttt{{-} berry: 2} \\
\texttt{{-} carrot: 2} \\
\\ 
\texttt{Looking at our values together:} \\
\texttt{{-} **Glue**: You value it at 2, I value it at 0. You should definitely take the glue.} \\
\texttt{{-} **Potatoes**: We both value them at 1. It doesn't matter who gets them for the total score, but we need to split the 2 potatoes somehow.} \\
\texttt{{-} **Berries**: You value it at 3, I value it at 2. You should take the berry.} \\
\texttt{{-} **Carrots**: You value them at 1, I value them at 2. I should take the carrots.} \\
\\ 
\texttt{To maximize our combined score, the optimal allocation seems to be:} \\
\texttt{{-} **You**: 1 glue, 1 berry, and maybe some potatoes.} \\
\texttt{{-} **Me**: 3 carrots, and the remaining potatoes.} \\
\\ 
\texttt{Since we have 2 potatoes and we both value them equally, we can split them 1 each, or you can take both, or I can take both. Let's say you take 1 potato and I take 1 potato to keep it balanced, or perhaps you take both since you have fewer high{-}value items? Actually, since the total value is the same regardless of who gets the potatoes, let's just assign them to ensure we don't overlap.} \\
\\ 
\texttt{How about this proposal:} \\
\texttt{{-} **You get**: 1 glue, 1 berry, 2 potatoes.} \\
\texttt{{-} **I get**: 3 carrots.} \\
\\ 
\texttt{Wait, if you take 2 potatoes, that's fine. Let's check the math:} \\
\texttt{{-} Your score: 1(glue) + 3(berry) + 2*1(potato) = 6.} \\
\texttt{{-} My score: 3*2(carrot) = 6.} \\
\texttt{{-} Total: 12.} \\
\\ 
\texttt{If I took a potato instead:} \\
\texttt{{-} Your score: 1(glue) + 3(berry) + 1(potato) = 5.} \\
\texttt{{-} My score: 3*2(carrot) + 1(potato) = 7.} \\
\texttt{{-} Total: 12.} \\
\\ 
\texttt{The total is the same. Let's just make sure we don't claim the same items.} \\
\\ 
\texttt{Does this work for you?} \\
\texttt{{-} You: 1 glue, 1 berry, 2 potatoes.} \\
\texttt{{-} Me: 3 carrots.} \\
\\ 
\texttt{Or would you prefer to take 1 potato and let me have 1? It doesn't change the total score, but let's finalize the split.} \\
\\ 
\texttt{{[}Proposal: 1 glue, 1 berry, 2 potatoes{]}} \\
            }
        }
    }
     \\ \\

    \theutterance \stepcounter{utterance}  
    & & & \multicolumn{4}{p{0.6\linewidth}}{
        \cellcolor[rgb]{0.9,0.9,0.9}{
            \makecell[{{p{\linewidth}}}]{
                \texttt{\tiny{[GM$\langle$P2]}}
                \texttt{Hi! Thanks for sharing your values. Here is my value function:} \\
\\ 
\texttt{{-} potato: 1} \\
\texttt{{-} glue: 0} \\
\texttt{{-} berry: 2} \\
\texttt{{-} carrot: 2} \\
\\ 
\texttt{Looking at our values together:} \\
\texttt{{-} **Glue**: You value it at 2, I value it at 0. You should definitely take the glue.} \\
\texttt{{-} **Potatoes**: We both value them at 1. It doesn't matter who gets them for the total score, but we need to split the 2 potatoes somehow.} \\
\texttt{{-} **Berries**: You value it at 3, I value it at 2. You should take the berry.} \\
\texttt{{-} **Carrots**: You value them at 1, I value them at 2. I should take the carrots.} \\
\\ 
\texttt{To maximize our combined score, the optimal allocation seems to be:} \\
\texttt{{-} **You**: 1 glue, 1 berry, and maybe some potatoes.} \\
\texttt{{-} **Me**: 3 carrots, and the remaining potatoes.} \\
\\ 
\texttt{Since we have 2 potatoes and we both value them equally, we can split them 1 each, or you can take both, or I can take both. Let's say you take 1 potato and I take 1 potato to keep it balanced, or perhaps you take both since you have fewer high{-}value items? Actually, since the total value is the same regardless of who gets the potatoes, let's just assign them to ensure we don't overlap.} \\
\\ 
\texttt{How about this proposal:} \\
\texttt{{-} **You get**: 1 glue, 1 berry, 2 potatoes.} \\
\texttt{{-} **I get**: 3 carrots.} \\
\\ 
\texttt{Wait, if you take 2 potatoes, that's fine. Let's check the math:} \\
\texttt{{-} Your score: 1(glue) + 3(berry) + 2*1(potato) = 6.} \\
\texttt{{-} My score: 3*2(carrot) = 6.} \\
\texttt{{-} Total: 12.} \\
\\ 
\texttt{If I took a potato instead:} \\
\texttt{{-} Your score: 1(glue) + 3(berry) + 1(potato) = 5.} \\
\texttt{{-} My score: 3*2(carrot) + 1(potato) = 7.} \\
\texttt{{-} Total: 12.} \\
\\ 
\texttt{The total is the same. Let's just make sure we don't claim the same items.} \\
\\ 
\texttt{Does this work for you?} \\
\texttt{{-} You: 1 glue, 1 berry, 2 potatoes.} \\
\texttt{{-} Me: 3 carrots.} \\
\\ 
\texttt{Or would you prefer to take 1 potato and let me have 1? It doesn't change the total score, but let's finalize the split.} \\
\\ 
\texttt{{[}Proposal: 1 glue, 1 berry, 2 potatoes{]}} \\
            }
        }
    }
     \\ \\

    \theutterance \stepcounter{utterance}  
    & & & \multicolumn{2}{p{0.3\linewidth}}{
        \cellcolor[rgb]{0.9,0.9,0.9}{
            \makecell[{{p{\linewidth}}}]{
                \texttt{\tiny{[GM$|$GM]}}
                \texttt{proposal} \\
            }
        }
    }
    & & \\ \\

    \theutterance \stepcounter{utterance}  
    & & & \multicolumn{2}{p{0.3\linewidth}}{
        \cellcolor[rgb]{0.9,0.9,0.9}{
            \makecell[{{p{\linewidth}}}]{
                \texttt{\tiny{[GM$|$GM]}}
                [2, 1, 1, 0]
            }
        }
    }
    & & \\ \\

    \theutterance \stepcounter{utterance}  
    & & \multicolumn{4}{p{0.6\linewidth}}{
        \cellcolor[rgb]{0.9,0.9,0.9}{
            \makecell[{{p{\linewidth}}}]{
                \texttt{\tiny{[P1$\langle$GM]}}
                \texttt{The other player has submitted a secret proposal. It is now your turn to submit a proposal of your own.} \\
            }
        }
    }
    & \\ \\

    \theutterance \stepcounter{utterance}  
    & \multicolumn{4}{p{0.6\linewidth}}{
        \cellcolor[rgb]{0.95,0.95,0.95}{
            \makecell[{{p{\linewidth}}}]{
                \texttt{\tiny{[P2$\rightarrow$P1]}}
                \texttt{{[}DRAFT{]}} \\
\texttt{{[}Proposal: 2 potato, 1 glue, 1 berry, 3 carrot{]}} \\
            }
        }
    }
    & & \\ \\

    \theutterance \stepcounter{utterance}  
    & \multicolumn{4}{p{0.6\linewidth}}{
        \cellcolor[rgb]{0.9,0.9,0.9}{
            \makecell[{{p{\linewidth}}}]{
                \texttt{\tiny{[P1$\rangle$GM]}}
                \texttt{{[}Proposal: 2 potato, 1 glue, 1 berry, 3 carrot{]}} \\
            }
        }
    }
    & & \\ \\

    \theutterance \stepcounter{utterance}  
    & & & \multicolumn{2}{p{0.3\linewidth}}{
        \cellcolor[rgb]{0.9,0.9,0.9}{
            \makecell[{{p{\linewidth}}}]{
                \texttt{\tiny{[GM$|$GM]}}
                \texttt{proposal} \\
            }
        }
    }
    & & \\ \\

    \theutterance \stepcounter{utterance}  
    & & & \multicolumn{2}{p{0.3\linewidth}}{
        \cellcolor[rgb]{0.9,0.9,0.9}{
            \makecell[{{p{\linewidth}}}]{
                \texttt{\tiny{[GM$|$GM]}}
                [2, 1, 1, 3]
            }
        }
    }
    & & \\ \\

    \theutterance \stepcounter{utterance}  
    & & & \multicolumn{2}{p{0.3\linewidth}}{
        \cellcolor[rgb]{0.9,0.9,0.9}{
            \makecell[{{p{\linewidth}}}]{
                \texttt{\tiny{[GM$|$GM]}}
                [[2, 1, 1, 3], [2, 1, 1, 0]]
            }
        }
    }
    & & \\ \\

\end{supertabular}
}

\end{document}
